\section{Conclusions}
\label{sec:conclusions}

We presented \bsi\ 2.0, an integrated platform for sequence analysis and protein
structure retrieval whose distinguishing feature is that every retrieved structure
is checked against the query and reported on an explicit five-level scale, with a
qualifier for confirmed but partial structures. On \N{ds-n} randomly sampled
proteins, judged against independently curated cross-references, accepting the top
sequence-search hit returned a structure not catalogued under the queried accession
for \N{n1-all-false} queries, including \N{n1-has-false} of \N{n1-has-acc} hits for
proteins that have a structure; accepting only M4 accepted \N{s1-all-false} wrong
structures, at the cost of returning a correct structure for \N{s1-has-sens}\% of
proteins that have one rather than \N{n1-has-sens}\%. The main reasons for that cost
are partial structures and entries that cannot be validated from a PDB-format file,
both addressable.

The same study qualifies the claim. Zero false associations at M4 is consistency
between two views of related curated data; the false associations that remain
at the permissive level are the same or a near-identical sequence; M2 is not
evidence of homology; and the prototype's accession discordances do not
demonstrate wrong structures. The software and the measurement are released
together, with the protocol amendments that record how the measurement was
corrected, so that the validation layer and the benchmark that tests it can both
be scrutinised.
